% Sec 3.3 insert -- REPLACES cc_channel_status_insert.tex from CC-PIPELINE-F.
% Targets channel_theory_outline.tex Section 3.3 (CC channel definition)
% and Section 9 (open problems list).

\subsection*{CC-channel anchor: the V$_{\text{quad}}$ recovery via Newton-polygon + Birkhoff (status post CC-PIPELINE-G)}

For a degree-2 partial continued fraction with linear $a_n=\alpha_1 n+\alpha_0$
and quadratic $b_n=\beta_2 n^2+\beta_1 n+\beta_0$, the partial-denominator
generating function $f(z)=\sum_{n\ge 0}Q_n z^n$ satisfies the order-2 linear
ODE
\begin{equation}
\Bigl[1
- z\bigl(\beta_2(\theta+1)^2+\beta_1(\theta+1)+\beta_0\bigr)
- z^2\bigl(\alpha_1(\theta+2)+\alpha_0\bigr)\Bigr]\,f(z)=1,
\qquad \theta=z\,\partial_z. \label{eq:Lf-cc}
\end{equation}
The Newton polygon of the homogeneous part of \eqref{eq:Lf-cc} at $z=0$ has a
single slope-$1/2$ edge of multiplicity 2, with characteristic polynomial
$1-(\beta_2/4)c^2$ and roots
\(c=\pm 2/\sqrt{\beta_2}\); equivalently, after the substitution $z=u^2$ the
operator carries a Gevrey-1-in-$u$ formal pair
\[
f_\pm(u)=\exp\bigl(\pm 2/(u\sqrt{\beta_2})\bigr)\,u^{\rho}\bigl(1+\textstyle\sum_{k\ge 1}a_k u^k\bigr),
\qquad \rho=-\tfrac32-\tfrac{\beta_1}{\beta_2}.
\]
This corrects the formal-series space $\mathcal D$ identified in
\textsc{CC-PIPELINE-F (UF-1)}: the literature constant $\xi_0=2/\sqrt{3}$ for
V$_{\text{quad}}$ is the Borel-singularity radius $|c|$ of \eqref{eq:Lf-cc}
(not of the BoT trans-series of $\log|\delta_n|$), and is recovered to
$200$ digits as the exact algebraic identity $\xi_0=2/\sqrt{\beta_2}|_{\beta_2=3}$.

\paragraph{Variant-A flip test in the corrected space.}
Apply the same pipeline to QL01, QL02, QL06, QL15, QL26 (PCF-1 v1.3 Sec~5).
Comparing the formal-solution invariants $(\xi_0,\rho,a_1,a_2,a_3,a_4)$
with the V$_{\text{quad}}$ P-III$(D_6)$ anchor at the
$15$-digit threshold:
$0/5$ flip (minimum agreement $<1$ digit on every non-V$_{\text{quad}}$ family).
The Borel-channel ``BOTH ARTEFACT'' verdict of \textsc{BOREL-CHANNEL-K-SCAN}
therefore stands, now with structural rather than asymptotic-tail evidence
(see Table~\ref{tab:cc-channel-v2}).

\paragraph{Cross-family Borel-radius coincidences.}
The radius $|c|=2/\sqrt{\beta_2}$ is determined entirely by the leading
coefficient $\beta_2$ of $b_n$.  Two pairs in our roster share radii:
$(\mathrm{V}_{\text{quad}},\mathrm{QL15})$ at $\beta_2=3$ and
$(\mathrm{QL01},\mathrm{QL02})$ at $\beta_2=1$.
The next invariant $\rho=-3/2-\beta_1/\beta_2$ separates V$_{\text{quad}}$
($\rho=-11/6$) from QL15 ($\rho=-5/6$) but coincides for the QL01/QL02
pair ($\rho=-5/2$ in both).  The first invariant that separates QL01 from
QL02 is $a_2$ (QL01: $17/512$; QL02: $-239/512$), which feeds from the
$z^2$ Newton-polygon points $(2,0),(2,1)$ encoding $a_n$.  The empirical
hierarchy of CC-channel invariants is therefore
\[
\xi_0~\succ~\rho~\succ~(a_1,a_2,\dots),
\]
with the $a_k$ tower carrying the $a_n$-dependence of the recurrence.

\paragraph{Status update for Section~9 (open problems).}
The previous \textbf{(op:cc-pipeline)} item is split:
\begin{itemize}
\item \textbf{(op:cc-formal-solution)} -- \emph{closed} by CC-PIPELINE-G:
the explicit Newton-polygon / Birkhoff pipeline above provides a
numerical CC datum at arbitrary precision for any degree-2 PCF.
\item \textbf{(op:cc-monodromy-RH)} -- \emph{open}: extend the formal-solution
extractor to a full Riemann-Hilbert monodromy datum at the irregular
singular point $z=0$ via Stokes-matrix computation.  This is the
prerequisite for a rigorous Painleve-class verdict at any
finite-digit threshold; the present session's verdict relies on the
heuristic fingerprint $(\xi_0,\rho,a_1,\dots,a_4)$ which is
necessary-not-sufficient.
\item \textbf{(op:cc-degree-d)} -- \emph{open}: generalise the Newton
polygon, char polynomial, and indicial formula to degree-$d$ PCFs
($a_n$ degree $d-1$, $b_n$ degree $d$).  Conjecture (preliminary):
$\xi_0=2\sqrt{(d-1)!/\beta_d}$ and $\rho=-(2d-1)/2-\beta_{d-1}/\beta_d$.
The $d=2$ case of this conjecture is established by CC-PIPELINE-G.
\end{itemize}

% End insert.  Recommended placement: after Definition~1 in Sec~3.3.
